% Компилируется xelatex (рекомендуется) или pdflatex
\documentclass[12pt,a4paper]{article}

\usepackage{iftex}
\ifXeTeX
  \usepackage{fontspec}
  \IfFontExistsTF{Times New Roman}{\setmainfont{Times New Roman}}{%
    \IfFontExistsTF{Liberation Serif}{\setmainfont{Liberation Serif}}{\setmainfont{DejaVu Serif}}}
  \usepackage[russian]{babel}
\else
  \usepackage[T2A]{fontenc}
  \usepackage[utf8]{inputenc}
  \usepackage[russian]{babel}
\fi

\usepackage[left=25mm,right=15mm,top=20mm,bottom=20mm]{geometry}
\usepackage{longtable}
\usepackage{array}
\usepackage{booktabs}
\usepackage{graphicx}
\usepackage{float}
\usepackage{enumitem}
\setlist{nosep}
\setlist[description]{leftmargin=0pt,itemsep=3pt}

\newcolumntype{L}[1]{>{\raggedright\arraybackslash}p{#1}}

\newenvironment{reqtable}{%
  \begin{longtable}{L{14mm} L{146mm}}
  \toprule
  \textbf{ID} & \textbf{Требование} \\
  \midrule
  \endhead
  \bottomrule
  \endfoot
}{\end{longtable}}

\newenvironment{uctable}[1]{%
  \par\medskip\noindent\begin{minipage}{\linewidth}
  \textbf{#1}\par\smallskip
  \renewcommand{\arraystretch}{1.25}%
  \begin{tabular}{@{}L{38mm} L{124mm}@{}}
  \toprule
}{%
  \bottomrule
  \end{tabular}
  \end{minipage}\par\medskip
}

\title{Информационная система «Бдительность»\\[4pt]
\large Этап 1: описание предметной области и требования}
\author{}
\date{}

\begin{document}
\maketitle
\tableofcontents
\newpage

%=====================================================================
\section{Описание предметной области}

\subsection{Общая характеристика}

Предметная область основана на сеттинге игры \emph{Beholder}: тоталитарное
государство, в котором власть контролирует граждан через жилой фонд.
Рассматривается один многоквартирный дом. За дом отвечает управляющий,
назначенный государством. Он обслуживает жильцов и одновременно следит за
их благонадёжностью.

Отношения жильца с государством строятся на двух механизмах. Первый ---
доносы: любой жилец может сообщить о нарушении, допущенном соседом, а
управляющий проверяет сообщение и назначает наказание. Второй ---
государственные услуги: всё, что жилец получает от государства, оформляется
заявкой, и доступ к услугам зависит от благонадёжности жильца.

Благонадёжность выражается числом --- индексом лояльности. Нарушения
снижают его и закрывают доступ к услугам, а полезное для государства
поведение повышает его и открывает новые услуги. Таким образом, доносы и
услуги образуют единый цикл.

Система реализуется в виде веб-сайта, оформленного в стиле портала
государственных услуг.

\subsection{Участники}

\begin{description}
  \item[Жилец] --- гражданин, проживающий в квартире дома. Пополняет
    баланс, оплачивает квартплату, подаёт заявки на услуги и доносы на
    соседей.
  \item[Управляющий] --- должностное лицо, ответственное за дом. Заселяет
    и выселяет жильцов, рассматривает и исполняет заявки, рассматривает
    доносы и назначает наказания.
\end{description}

\subsection{Основные сущности и термины}

\begin{description}
  \item[Дом] --- объект учёта, состоящий из квартир. Система обслуживает
    один дом.
  \item[Квартира] --- жилое помещение, за которое начисляется квартплата.
    Квартира может быть свободна или занята.
  \item[Анкета] --- учётные данные жильца.
  \item[Код приглашения] --- одноразовый код, выдаваемый управляющим при
    заселении для активации учётной записи. Самостоятельная регистрация
    невозможна.
  \item[Индекс лояльности (ИЛ)] --- число, отражающее благонадёжность
    жильца. Изменяется автоматически по фиксированным правилам и определяет
    статус благонадёжности, набор доступных услуг и очерёдность
    рассмотрения заявок.
  \item[Статус благонадёжности] --- категория жильца, определяемая
    значением ИЛ.
  \item[Коммунальный баланс] --- внутренний счёт жильца. Пополняется через
    платёжный сервис ЮKassa, расходуется на квартплату, пошлины и штрафы.
  \item[Квартплата] --- ежемесячное начисление за квартиру. Просрочка
    оплаты является нарушением.
  \item[Услуга] --- предоставляемая государством услуга. Перечень услуг
    задан в системе. Для каждой услуги определены требуемый ИЛ, пошлина и
    форма результата.
  \item[Заявка] --- обращение жильца по услуге. Имеет статус, историю
    обработки и результат.
  \item[Результат заявки] --- итог исполнения: заключение управляющего,
    выданный документ или разрешение.
  \item[Донос] --- сообщение жильца о нарушении, допущенном другим жильцом.
    Может быть анонимным.
  \item[Наказание] --- решение управляющего по подтверждённому доносу:
    предупреждение, штраф или выселение. Любое наказание снижает ИЛ
    нарушителя.
  \item[Досье] --- накопительное дело на жильца: доносы, наказания,
    история ИЛ, заметки управляющего. Жильцу недоступно.
\end{description}

\subsection{Перечень услуг}

Перечень услуг фиксирован. Услуги упорядочены по возрастанию требуемого ИЛ:
чем ценнее услуга, тем выше должна быть благонадёжность жильца.

\begin{longtable}{L{56mm} L{12mm} L{24mm} L{58mm}}
\toprule
\textbf{Услуга} & \textbf{ИЛ} & \textbf{Пошлина} & \textbf{Результат} \\
\midrule
\endhead
\bottomrule
\endfoot
Бытовой ремонт & 0 & нет & Отметка о выполнении работ \\
Справка о проживании & 20 & низкая & Документ \\
Продовольственный талон & 30 & низкая & Талон со сроком действия \\
Разрешение на приём гостей & 40 & средняя & Разрешение на указанные даты \\
Разрешение на выход в комендантский час & 50 & средняя & Разрешение на указанную дату \\
Замена мебели и бытовой техники & 60 & высокая & Отметка о выполнении работ \\
Переселение в другую квартиру & 70 & высокая & Смена квартиры \\
Разрешение на выезд из города & 80 & высокая & Разрешение на период \\
Снятие предупреждения из досье & 85 & очень высокая & Предупреждение снято \\
\end{longtable}

\subsection{Основные процессы}

\paragraph{Заселение.} Управляющий заселяет жильца в свободную квартиру и
выдаёт ему код приглашения. Жилец активирует учётную запись и получает
начальный ИЛ.

\paragraph{Платежи.} Жилец пополняет баланс через ЮKassa. Каждый месяц
система начисляет квартплату, которую жилец оплачивает с баланса. Просрочка
оплаты снижает ИЛ.

\paragraph{Заявка на услугу.} Жилец выбирает доступную ему услугу и подаёт
заявку, оплачивая пошлину. Управляющий рассматривает заявку, при
необходимости запрашивает уточнение, затем исполняет её с оформлением
результата либо отклоняет с указанием причины. Жилец следит за статусом
заявки и получает результат в личном кабинете.

\paragraph{Донос.} Жилец подаёт донос на соседа. Управляющий рассматривает
его и либо подтверждает с назначением наказания, либо отклоняет, либо
признаёт ложным. Подтверждённый донос повышает ИЛ доносчика, ложный ---
снижает. Управляющий также может наказать жильца без доноса.

\paragraph{Выселение.} При падении ИЛ ниже критического уровня жилец
подлежит выселению. Управляющий выселяет жильца, его учётная запись
блокируется, квартира освобождается.

%=====================================================================
\section{Назначение информационной системы}

Без информационной системы обращения жильцов теряются, сроки их
рассмотрения не контролируются, доносы не проверяются, а решения
управляющего невозможно отследить. Система позволяет:

\begin{enumerate}
  \item вести учёт квартир и жильцов дома;
  \item предоставлять государственные услуги в электронном виде и
        отслеживать обработку заявок;
  \item ставить доступ к услугам в зависимость от благонадёжности жильца;
  \item принимать и рассматривать доносы, назначать наказания;
  \item автоматически оценивать благонадёжность жильцов через индекс
        лояльности;
  \item вести коммунальный баланс жильцов: пополнение, квартплату, пошлины
        и штрафы;
  \item накапливать сведения о жильцах в досье.
\end{enumerate}

%=====================================================================
\section{Функциональные требования}

Требования сгруппированы по акторам: жилец (U), управляющий (C), система (S).

\subsection{Жилец}

\subsubsection*{Учётная запись и профиль}
\begin{reqtable}
U001 & Жилец должен иметь возможность активировать учётную запись по коду приглашения, полученному от управляющего, и задать пароль. Код используется однократно. \\
U002 & Жилец должен иметь возможность войти в систему по логину и паролю и выйти из неё. Выселенный жилец войти в систему не может. \\
U003 & Жилец должен иметь возможность просматривать свою анкету и сведения о закреплённой за ним квартире. \\
U004 & Жилец должен иметь возможность просматривать текущее значение своего ИЛ, статус благонадёжности и историю изменений ИЛ с указанием даты и причины каждого изменения. Если изменение вызвано доносом, автор доноса не указывается. \\
\end{reqtable}

\subsubsection*{Баланс и платежи}
\begin{reqtable}
U005 & Жилец должен иметь возможность пополнить коммунальный баланс на выбранную сумму через ЮKassa. После оплаты жилец возвращается в личный кабинет и видит результат операции. \\
U006 & Жилец должен иметь возможность просматривать текущий баланс, наличие задолженности и историю операций: пополнения, квартплату, пошлины, штрафы и возвраты. \\
U007 & Жилец должен иметь возможность просматривать начисления квартплаты и оплачивать их с баланса. \\
\end{reqtable}

\subsubsection*{Заявки на услуги}
\begin{reqtable}
U008 & Жилец должен иметь возможность просматривать перечень услуг с описанием, требуемым ИЛ, размером пошлины и формой результата. Услуги, недоступные по текущему ИЛ, отображаются, но подать заявку на них нельзя. \\
U009 & Жилец должен иметь возможность подать заявку на доступную ему услугу, указав необходимые сведения. При подаче с баланса списывается пошлина. Подать заявку при недостатке средств или наличии задолженности нельзя. \\
U010 & Жилец должен иметь возможность отменить свою заявку, пока управляющий не принял её в работу. При отмене пошлина возвращается на баланс. \\
U011 & Жилец должен иметь возможность просматривать список своих заявок и карточку каждой заявки: текущий статус, историю изменения статусов и комментарии управляющего. \\
U012 & Жилец должен иметь возможность ответить на запрос уточнения по заявке, после чего заявка возвращается управляющему на рассмотрение. \\
U013 & Жилец должен иметь возможность получить результат выполненной заявки: заключение управляющего и выданный документ или разрешение. При отклонении заявки жилец видит причину отказа. \\
U014 & Жилец должен иметь возможность просматривать полученные разрешения и документы, а также срок их действия. \\
\end{reqtable}

\subsubsection*{Доносы и наказания}
\begin{reqtable}
U015 & Жилец должен иметь возможность подать донос на другого жильца дома, указав суть нарушения и обстоятельства. Донос на самого себя подать нельзя. \\
U016 & Жилец должен иметь возможность подать донос анонимно. В этом случае его личность не раскрывается ни управляющему, ни нарушителю, а ИЛ доносчика по итогам рассмотрения не изменяется. \\
U017 & Жилец должен иметь возможность просматривать свои доносы и итог их рассмотрения: подтверждён, отклонён или признан ложным. \\
U018 & Жилец должен иметь возможность просматривать назначенные ему наказания с описанием нарушения и последствиями для баланса и ИЛ. \\
U019 & Жилец должен получать в личном кабинете уведомления об изменении статуса заявок, итогах рассмотрения доносов, начислениях, штрафах и изменениях ИЛ. \\
\end{reqtable}

\subsection{Управляющий}

\subsubsection*{Жилой фонд и жильцы}
\begin{reqtable}
C001 & Управляющий должен иметь возможность просматривать квартиры дома с указанием занятости и проживающих жильцов. \\
C002 & Управляющий должен иметь возможность заселить нового жильца в свободную квартиру, заполнив его анкету. При заселении система выдаёт код приглашения для активации учётной записи. \\
C003 & Управляющий должен иметь возможность просматривать список жильцов с их ИЛ, статусом благонадёжности и состоянием баланса. \\
C004 & Управляющий должен иметь возможность изменять анкетные данные жильцов. \\
C005 & Управляющий должен иметь возможность выселить жильца. При выселении учётная запись жильца блокируется, а квартира становится свободной. \\
\end{reqtable}

\subsubsection*{Обработка заявок}
\begin{reqtable}
C006 & Управляющий должен иметь возможность просматривать очередь поступивших заявок. Заявки жильцов с более высоким ИЛ находятся выше в очереди. \\
C007 & Управляющий должен иметь возможность открыть заявку и просмотреть её содержание, сведения о заявителе и историю обработки. \\
C008 & Управляющий должен иметь возможность принять заявку в работу. После этого жилец не может её отменить. \\
C009 & Управляющий должен иметь возможность запросить у жильца уточнение по заявке с комментарием. \\
C010 & Управляющий должен иметь возможность отклонить заявку с обязательным указанием причины и решить, возвращается ли жильцу пошлина. \\
C011 & Управляющий должен иметь возможность завершить заявку, оформив результат: заключение и, если это предусмотрено услугой, документ или разрешение со сроком действия. \\
\end{reqtable}

\subsubsection*{Доносы и наказания}
\begin{reqtable}
C012 & Управляющий должен иметь возможность просматривать поступившие доносы. Автор анонимного доноса не отображается. \\
C013 & Управляющий должен иметь возможность подтвердить донос и назначить нарушителю наказание: предупреждение, штраф с указанием суммы или выселение. \\
C014 & Управляющий должен иметь возможность отклонить донос за недостатком оснований. В этом случае ИЛ участников не изменяется. \\
C015 & Управляющий должен иметь возможность признать донос ложным. В этом случае снижается ИЛ доносчика. \\
C016 & Управляющий должен иметь возможность самостоятельно зафиксировать нарушение жильца и назначить ему наказание без доноса. \\
C017 & Управляющий должен иметь возможность просматривать досье жильца: доносы на него, назначенные наказания, историю ИЛ, заявки и платежи. \\
C018 & Управляющий должен иметь возможность добавлять в досье жильца собственные заметки. \\
\end{reqtable}

\subsection{Система}
\begin{reqtable}
S001 & Система должна ежемесячно начислять квартплату жильцам каждой занятой квартиры. \\
S002 & Система должна фиксировать нарушение при просрочке оплаты квартплаты и снижать ИЛ жильца. \\
S003 & Система должна зачислять средства на баланс только после подтверждения оплаты от ЮKassa, ровно один раз на каждый платёж. \\
S004 & Система должна списывать назначенный штраф с баланса жильца. При недостатке средств образуется задолженность, которая погашается при следующем пополнении. \\
S005 & Система должна автоматически пересчитывать ИЛ по фиксированным правилам при наказании, подтверждении или признании ложным доноса, своевременной оплате и просрочке квартплаты. \\
S006 & Система должна автоматически обновлять статус благонадёжности и набор доступных жильцу услуг при изменении ИЛ. \\
S007 & Система должна уведомлять управляющего о жильцах, чей ИЛ опустился ниже критического уровня и которые подлежат выселению. \\
S008 & Система должна вести историю всех операций с балансом и изменений ИЛ, которую нельзя изменить или удалить. \\
\end{reqtable}

%=====================================================================
\section{Нефункциональные требования}

\begin{reqtable}
N001 & \textbf{Платформа.} Веб-приложение на React, Spring и PostgreSQL, развёртываемое на сервере helios. \\
N002 & \textbf{Платежи.} Пополнение баланса выполняется через ЮKassa в тестовом режиме без использования реальных средств. \\
N003 & \textbf{Производительность.} Типовые операции выполняются без заметных для пользователя задержек при одновременной работе всех жильцов дома. \\
N004 & \textbf{Безопасность.} Пароли хранятся в защищённом виде, права доступа проверяются на сервере, приложение защищено от распространённых веб-атак. \\
N005 & \textbf{Конфиденциальность.} Автор анонимного доноса не раскрывается. \\
N006 & \textbf{Целостность.} Финансовые операции, назначение наказаний и пересчёт ИЛ выполняются атомарно. \\
N007 & \textbf{Надёжность.} Ошибки обрабатываются корректно, пользователь получает понятные сообщения; повторное подтверждение оплаты не приводит к двойному зачислению. \\
N008 & \textbf{Интерфейс.} Единый строгий стиль портала государственных услуг, интерфейс на русском языке. \\
N009 & \textbf{Совместимость.} Корректная работа в современных браузерах на компьютерах и мобильных устройствах. \\
N010 & \textbf{Сопровождаемость.} Многоуровневая архитектура с разделением представления, бизнес-логики и доступа к данным. \\
\end{reqtable}

%=====================================================================
\section{Прецеденты}

Ниже приведено описание прецедентов системы. Номер шага в альтернативном
потоке указывает на шаг основного потока, от которого он ответвляется.

\subsubsection*{Прецеденты актора «Жилец»}
\begin{uctable}{UC01. Активировать учётную запись}
Актор & Жилец \\
Цель & Получить доступ к личному кабинету. \\
Предусловия & Управляющий заселил жильца и выдал ему код приглашения. \\
Основной поток & 1.~Жилец открывает страницу активации. \newline 2.~Жилец вводит код приглашения. \newline 3.~Система проверяет код. \newline 4.~Жилец задаёт логин и пароль. \newline 5.~Система создаёт учётную запись, назначает начальный ИЛ и открывает личный кабинет. \\
Альтернативные потоки & 3а. Код неверен или уже использован: система сообщает об ошибке, активация не выполняется. \\
Постусловия & Учётная запись активна, код приглашения больше не действует. \\
\end{uctable}

\begin{uctable}{UC02. Войти в личный кабинет}
Актор & Жилец \\
Цель & Начать работу с системой. \\
Предусловия & Учётная запись активирована. \\
Основной поток & 1.~Жилец вводит логин и пароль. \newline 2.~Система проверяет данные и статус жильца. \newline 3.~Система открывает личный кабинет. \\
Альтернативные потоки & 2а. Данные неверны: система сообщает об ошибке. \newline 2б. Жилец выселен: вход запрещён. \\
Постусловия & Жилец авторизован в системе. \\
\end{uctable}

\begin{uctable}{UC03. Просмотреть перечень услуг}
Актор & Жилец \\
Цель & Узнать, какие услуги доступны жильцу. \\
Предусловия & Жилец авторизован. \\
Основной поток & 1.~Жилец открывает раздел услуг. \newline 2.~Система показывает перечень услуг с требуемым ИЛ, пошлиной и формой результата. \newline 3.~Система отмечает услуги, недоступные по текущему ИЛ жильца. \\
Альтернативные потоки & — \\
Постусловия & — \\
\end{uctable}

\begin{uctable}{UC04. Пополнить баланс}
Актор & Жилец \\
Цель & Внести средства на коммунальный баланс. \\
Предусловия & Жилец авторизован. \\
Основной поток & 1.~Жилец указывает сумму пополнения. \newline 2.~Система создаёт платёж в ЮKassa и перенаправляет жильца на страницу оплаты. \newline 3.~Жилец оплачивает платёж. \newline 4.~ЮKassa подтверждает оплату. \newline 5.~Система зачисляет средства на баланс; если есть задолженность, она погашается в первую очередь. \newline 6.~Жилец возвращается в личный кабинет и видит результат. \\
Альтернативные потоки & 3а. Жилец отменяет оплату: платёж отменяется, баланс не изменяется. \\
Постусловия & Баланс увеличен, операция отражена в истории. \\
\end{uctable}

\begin{uctable}{UC05. Просмотреть баланс и историю операций}
Актор & Жилец \\
Цель & Проконтролировать состояние счёта. \\
Предусловия & Жилец авторизован. \\
Основной поток & 1.~Жилец открывает раздел баланса. \newline 2.~Система показывает текущий баланс, задолженность и список операций. \\
Альтернативные потоки & — \\
Постусловия & — \\
\end{uctable}

\begin{uctable}{UC06. Оплатить квартплату}
Актор & Жилец \\
Цель & Погасить начисление за квартиру. \\
Предусловия & Система начислила квартплату. \\
Основной поток & 1.~Жилец открывает список начислений. \newline 2.~Жилец выбирает неоплаченное начисление и подтверждает оплату. \newline 3.~Система списывает сумму с баланса и отмечает начисление оплаченным. \newline 4.~Система пересчитывает ИЛ за своевременную оплату. \\
Альтернативные потоки & 3а. Недостаточно средств: система предлагает пополнить баланс. \\
Постусловия & Начисление оплачено. \\
\end{uctable}

\begin{uctable}{UC07. Подать заявку на услугу}
Актор & Жилец \\
Цель & Получить государственную услугу. \\
Предусловия & Жилец авторизован, услуга доступна по его ИЛ. \\
Основной поток & 1.~Жилец выбирает услугу в перечне. \newline 2.~Жилец заполняет сведения по заявке. \newline 3.~Система проверяет ИЛ, баланс и отсутствие задолженности. \newline 4.~Система списывает пошлину и регистрирует заявку в статусе «Подана». \newline 5.~Заявка попадает в очередь управляющего. \\
Альтернативные потоки & 3а. ИЛ ниже требуемого: подача невозможна. \newline 3б. Недостаточно средств или есть задолженность: система предлагает пополнить баланс. \\
Постусловия & Заявка зарегистрирована, пошлина списана. \\
\end{uctable}

\begin{uctable}{UC08. Отменить заявку}
Актор & Жилец \\
Цель & Отказаться от поданной заявки. \\
Предусловия & Заявка находится в статусе «Подана». \\
Основной поток & 1.~Жилец открывает заявку и выбирает отмену. \newline 2.~Система переводит заявку в статус «Отменена». \newline 3.~Система возвращает пошлину на баланс. \\
Альтернативные потоки & 1а. Заявка уже принята в работу: отмена недоступна. \\
Постусловия & Заявка отменена, пошлина возвращена. \\
\end{uctable}

\begin{uctable}{UC09. Ответить на запрос уточнения}
Актор & Жилец \\
Цель & Дополнить заявку по требованию управляющего. \\
Предусловия & Заявка находится в статусе «Требует уточнения». \\
Основной поток & 1.~Жилец получает уведомление и открывает заявку. \newline 2.~Жилец читает комментарий управляющего и вводит ответ. \newline 3.~Система возвращает заявку в статус «В работе». \\
Альтернативные потоки & — \\
Постусловия & Заявка снова на рассмотрении у управляющего. \\
\end{uctable}

\begin{uctable}{UC10. Получить результат заявки}
Актор & Жилец \\
Цель & Узнать итог рассмотрения заявки. \\
Предусловия & Заявка завершена или отклонена. \\
Основной поток & 1.~Жилец получает уведомление и открывает заявку. \newline 2.~Система показывает заключение управляющего и выданный документ или разрешение. \newline 3.~Жилец при необходимости скачивает документ. \\
Альтернативные потоки & 2а. Заявка отклонена: система показывает причину отказа и сведения о возврате пошлины. \\
Постусловия & — \\
\end{uctable}

\begin{uctable}{UC11. Подать донос}
Актор & Жилец \\
Цель & Сообщить о нарушении, допущенном соседом. \\
Предусловия & Жилец авторизован. \\
Основной поток & 1.~Жилец выбирает жильца, на которого подаёт донос. \newline 2.~Жилец описывает нарушение и обстоятельства. \newline 3.~Жилец указывает, подаётся ли донос анонимно. \newline 4.~Система регистрирует донос и передаёт его управляющему. \newline 5.~Жилец может отслеживать статус доноса. \\
Альтернативные потоки & 1а. Жилец выбрал самого себя: система не допускает подачу. \\
Постусловия & Донос зарегистрирован и ожидает рассмотрения. \\
\end{uctable}

\begin{uctable}{UC12. Просмотреть историю ИЛ}
Актор & Жилец \\
Цель & Узнать свой уровень благонадёжности и причины его изменения. \\
Предусловия & Жилец авторизован. \\
Основной поток & 1.~Жилец открывает раздел ИЛ. \newline 2.~Система показывает текущий ИЛ, статус благонадёжности и историю изменений с датами и причинами. \\
Альтернативные потоки & 2а. Изменение вызвано доносом: автор доноса не указывается. \\
Постусловия & — \\
\end{uctable}

\subsubsection*{Прецеденты актора «Управляющий»}
\begin{uctable}{UC13. Заселить жильца}
Актор & Управляющий \\
Цель & Зарегистрировать нового жильца дома. \\
Предусловия & В доме есть свободная квартира. \\
Основной поток & 1.~Управляющий выбирает свободную квартиру. \newline 2.~Управляющий заполняет анкету жильца. \newline 3.~Система создаёт запись о жильце, досье и код приглашения. \newline 4.~Управляющий передаёт код жильцу. \\
Альтернативные потоки & 2а. Анкета заполнена не полностью: система требует заполнить обязательные поля. \\
Постусловия & Квартира занята, жилец может активировать учётную запись. \\
\end{uctable}

\begin{uctable}{UC14. Выселить жильца}
Актор & Управляющий \\
Цель & Освободить квартиру и прекратить доступ жильца к системе. \\
Предусловия & Жилец проживает в доме. \\
Основной поток & 1.~Управляющий открывает карточку жильца и выбирает выселение. \newline 2.~Управляющий указывает основание. \newline 3.~Система блокирует учётную запись, освобождает квартиру и фиксирует выселение в досье. \\
Альтернативные потоки & 1а. Выселение инициировано системой при падении ИЛ ниже критического уровня: управляющий получает уведомление и подтверждает выселение. \\
Постусловия & Жилец выселен, квартира свободна. \\
\end{uctable}

\begin{uctable}{UC15. Посмотреть досье}
Актор & Управляющий \\
Цель & Изучить сведения о жильце. \\
Предусловия & Жилец зарегистрирован в системе. \\
Основной поток & 1.~Управляющий выбирает жильца. \newline 2.~Система показывает досье: доносы на жильца, наказания, историю ИЛ, заявки, платежи и заметки. \\
Альтернативные потоки & — \\
Постусловия & — \\
\end{uctable}

\begin{uctable}{UC16. Добавить заметку в досье}
Актор & Управляющий \\
Цель & Зафиксировать собственные наблюдения о жильце. \\
Предусловия & Управляющий открыл досье жильца. \\
Основной поток & 1.~Управляющий вводит текст заметки. \newline 2.~Система сохраняет заметку с датой. \\
Альтернативные потоки & — \\
Постусловия & Заметка добавлена в досье. \\
\end{uctable}

\begin{uctable}{UC17. Посмотреть список заявок}
Актор & Управляющий \\
Цель & Выбрать заявку для обработки. \\
Предусловия & Управляющий авторизован. \\
Основной поток & 1.~Управляющий открывает очередь заявок. \newline 2.~Система показывает заявки, упорядоченные по ИЛ заявителей. \newline 3.~Управляющий открывает карточку заявки с содержанием, сведениями о заявителе и историей обработки. \\
Альтернативные потоки & — \\
Постусловия & — \\
\end{uctable}

\begin{uctable}{UC18. Принять заявку в работу}
Актор & Управляющий \\
Цель & Начать исполнение заявки. \\
Предусловия & Заявка в статусе «Подана». \\
Основной поток & 1.~Управляющий открывает заявку и выбирает «Принять в работу». \newline 2.~Система переводит заявку в статус «В работе» и уведомляет жильца. \\
Альтернативные потоки & — \\
Постусловия & Заявка в работе, жилец больше не может её отменить. \\
\end{uctable}

\begin{uctable}{UC19. Запросить уточнение}
Актор & Управляющий \\
Цель & Получить от жильца недостающие сведения. \\
Предусловия & Заявка в статусе «В работе». \\
Основной поток & 1.~Управляющий вводит комментарий с вопросом. \newline 2.~Система переводит заявку в статус «Требует уточнения» и уведомляет жильца. \\
Альтернативные потоки & — \\
Постусловия & Заявка ожидает ответа жильца. \\
\end{uctable}

\begin{uctable}{UC20. Отклонить заявку}
Актор & Управляющий \\
Цель & Отказать в предоставлении услуги. \\
Предусловия & Заявка в статусе «Подана» или «В работе». \\
Основной поток & 1.~Управляющий выбирает отклонение и указывает причину. \newline 2.~Управляющий решает, возвращается ли пошлина. \newline 3.~Система переводит заявку в статус «Отклонена», при необходимости возвращает пошлину и уведомляет жильца. \\
Альтернативные потоки & 1а. Причина не указана: система не позволяет отклонить заявку. \\
Постусловия & Заявка отклонена. \\
\end{uctable}

\begin{uctable}{UC21. Завершить заявку}
Актор & Управляющий \\
Цель & Выдать жильцу результат услуги. \\
Предусловия & Заявка в статусе «В работе». \\
Основной поток & 1.~Управляющий оформляет заключение. \newline 2.~Если услуга предусматривает документ или разрешение, управляющий прикладывает его и указывает срок действия. \newline 3.~Система переводит заявку в статус «Выполнена» и уведомляет жильца. \newline 4.~Если услуга меняет данные жильца, например переселение, система применяет изменения. \\
Альтернативные потоки & — \\
Постусловия & Заявка выполнена, результат доступен жильцу. \\
\end{uctable}

\begin{uctable}{UC22. Рассмотреть донос}
Актор & Управляющий \\
Цель & Проверить сообщение о нарушении и принять решение. \\
Предусловия & Поступил донос. \\
Основной поток & 1.~Управляющий открывает донос и досье подозреваемого. \newline 2.~Управляющий оценивает обоснованность доноса. \newline 3.~Управляющий подтверждает донос и переходит к назначению наказания. \newline 4.~Система повышает ИЛ доносчика, если донос не анонимный. \\
Альтернативные потоки & 3а. Оснований недостаточно: донос отклоняется, ИЛ участников не изменяется. \newline 3б. Донос ложный: донос признаётся ложным, ИЛ доносчика снижается. \\
Постусловия & Донос рассмотрен, доносчик видит итог. \\
\end{uctable}

\begin{uctable}{UC23. Назначить наказание}
Актор & Управляющий \\
Цель & Наказать жильца за нарушение. \\
Предусловия & Донос подтверждён, либо управляющий самостоятельно выявил нарушение. \\
Основной поток & 1.~Управляющий выбирает вид наказания: предупреждение, штраф или выселение. \newline 2.~Для штрафа управляющий указывает сумму. \newline 3.~Система фиксирует наказание в досье, снижает ИЛ нарушителя и уведомляет его. \newline 4.~Система списывает штраф с баланса. \\
Альтернативные потоки & 1а. Выбрано выселение: выполняется прецедент «Выселить жильца». \newline 4а. Недостаточно средств: образуется задолженность. \\
Постусловия & Наказание назначено и отражено в досье. \\
\end{uctable}

%=====================================================================
\section{Используемые технологии}

Система строится по клиент-серверной архитектуре. Клиентская часть
представляет собой одностраничное веб-приложение, которое взаимодействует
с серверной частью через REST API. Серверная часть содержит бизнес-логику
и работает с базой данных.

\begin{longtable}{L{30mm} L{38mm} L{85mm}}
\toprule
\textbf{Компонент} & \textbf{Технология} & \textbf{Назначение} \\
\midrule
\endhead
\bottomrule
\endfoot
Клиентская часть & React & Пользовательский интерфейс личных кабинетов жильца и управляющего. \\
Серверная часть & Spring & Бизнес-логика, REST API, аутентификация и разграничение доступа, обработка уведомлений об оплате. \\
База данных & PostgreSQL & Хранение данных о квартирах, жильцах, заявках, доносах, наказаниях, платежах и истории ИЛ. \\
Платёжный сервис & ЮKassa (тестовый режим) & Пополнение коммунального баланса. \\
\end{longtable}

\end{document}